\section*{Hamiltonian identification theorem boundary}
\label{sec:hamiltonian-identification-boundary}

This appendix records the canonical audited boundary for the Yang--Mills-specific
identification step. The repository status is \emph{theorem boundary}, not a
completed proof of the four-dimensional $SU(3)$ mass-gap conjecture.

\subsection*{Three identification clauses already isolated}
The remaining bridge is organized around three abstract clauses that are already
separated cleanly in the audit:
\begin{itemize}
\item the \emph{vacuum clause}, identifying the convergence projector with the
vacuum-complement sector of the OS reconstruction;
\item the \emph{algebra clause}, identifying the convergence observable algebra
with the Wilson-loop / GNS algebra and its separating family;
\item the \emph{generator clause}, identifying the limiting semigroup generator
and the operative common core with the OS Hamiltonian domain.
\end{itemize}
In this reduced formulation the vacuum clause, algebra clause, and generator
clause are not the open problem by themselves; the open problem is the
Yang--Mills-specific realization needed to verify them in one common model.
The packet reductions below are the mechanism that tries to make that
realization finite and checkable. They use the dictionary lemma, the packet
generators, a common finite domain, and the OS word-expectation formula.

\subsection*{Finite packet reductions that are already abstractly available}
The current audit treats the following reductions as abstractly available and
therefore not the final Yang--Mills obstruction:
\begin{itemize}
\item finite Wilson-loop packet truncation is abstractly available;
\item finite rational-time packet truncation is abstractly available;
\item the finite Wilson-loop moment packet is the state-side finite datum;
\item a common-realization witness package is the bookkeeping device used to
compare the convergence and OS sides on finite packets.
\end{itemize}

\subsection*{Wilson-loop packet branch on $W_N$}
The state-side branch is the finite Wilson-loop packet entrywise equality theorem boundary.
Its internal reduction is the following theorem-sized ladder:
\begin{enumerate}
\item common-Schwinger-family criterion for Wilson-loop packet equality;
\item common-realization witness package for Wilson-loop packet equality;
\item tuple-fiber criterion for packet Schwinger witness existence;
\item common packet Schwinger witness from two descents;
\item the packet comparison factors through the tuple fibers;
\item common continuum Schwinger family theorem boundary on $W_N$;
\item convergence-side tuple-fiber constancy theorem boundary on $W_N$;
\item OS-side tuple-fiber constancy theorem boundary on $W_N$;
\item OS word formula implies tuple-fiber constancy;
\item tuple-image witness identification theorem boundary on $W_N$;
\item convergence-side Schwinger realization theorem boundary on $W_N$;
\item OS-side Schwinger realization theorem boundary on $W_N$.
\end{enumerate}
The shortcut architecture is also explicit: injective packet dictionary implies tuple-fiber constancy; letter-level injectivity implies packet-dictionary injectivity; and the letter-level injectivity theorem boundary on the packet alphabet is reduced to the statement that generator injectivity plus reflection-extension injectivity imply letter-level injectivity. That leaves the
generator-level injectivity theorem boundary on the finite packet generators,
the reflection-extension injectivity theorem boundary on the packet alphabet,
and therefore a shared shortcut obstruction. The intended shortcut is an injective single-letter dictionary,
and the scheduler now records this injectivity shortcut explicitly together with the unresolved shortcut subgoals.
This is why the appendix keeps both the exact live state-side blocker and the
blocked-action view visible at once: the dictionary shortcut is shared even when
the two packet branches are written separately.

\subsection*{Rational-time packet branch on $F_N$}
The action-side branch is the finite rational-time packet entrywise equality theorem boundary.
Its finite reduction is:
\begin{enumerate}
\item common Euclidean-action criterion for rational-time packet equality;
\item common-realization witness package for rational-time packet equality;
\item representative-fiber criterion for packet action witness existence;
\item common packet action witness from two descents;
\item packet representative fibers;
\item common Euclidean action theorem boundary on $F_N$;
\item convergence-side representative-fiber constancy theorem boundary on $F_N$;
\item OS-side representative-fiber constancy theorem boundary on $F_N$;
\item injective representative map implies representative-fiber constancy;
\item labeled representative encoding implies representative-map injectivity;
\item representative-image witness identification theorem boundary on $F_N$;
\item convergence-side action realization theorem boundary on $F_N$;
\item OS-side action realization theorem boundary on $F_N$.
\end{enumerate}
The same generator-level injectivity theorem boundary on the finite packet
generators enters here because the representative map must remain faithful on the
same packet alphabet that feeds the state-side dictionary.

\subsection*{Dictionary infrastructure shared by both branches}
The packet reductions only make sense after the following common dictionary
statements are isolated:
\begin{itemize}
\item single-letter dictionary theorem boundary;
\item adjoint/reflection dictionary theorem boundary;
\item word-to-tuple assembly theorem boundary;
\item tuple-domain admissibility theorem boundary.
\end{itemize}
These shared preliminaries explain why both packet branches are localized but not
independent: one common-realization witness package still depends on a single
Yang--Mills-specific identification mechanism.

\subsection*{Exact live theorem boundary}
The exact live state-side blocker is the convergence-side Schwinger realization
 theorem boundary on $W_N$: one must prove that the convergence-side continuum
state is realized by the same continuum Schwinger family used on the OS side.
The parallel action-side blocker is the convergence-side action realization
 theorem boundary on $F_N$: one must prove that the convergence-side packet
action is realized by the same Euclidean time-translation action that enters the
OS Hamiltonian construction. The current verifier now also requires the
Schwinger, OS reconstruction, operator-convergence, and semigroup artifacts to
carry the same action-specification hash before this action-side support is
accepted. The downstream countable localization layer uses explicit continuity
estimates: for contraction letters,
\[
\|g_1\cdots g_n-d_1\cdots d_n\|\leq n\varepsilon,
\]
so state differences are bounded by $2n\varepsilon$, and rational-time action
agreement is upgraded to real-time agreement by strong continuity. Until the
packet realization statements are proved,
packetwise equalities cannot be extended to the full continuum Wilson-loop state
or the full Euclidean semigroup.

Any manuscript claim that the mass gap is already proved must therefore first
discharge the vacuum clause, algebra clause, and generator clause in the same
Yang--Mills setting by proving the packet dictionary, the two packet realization
theorems, and the downstream countable extension steps.
\endinput
